\section{Summary}\label{sec:tsps-summary}

This chapter addresses the gap between structure-preserving signatures and
their threshold counterparts. Structure-preserving signatures integrate
naturally with non-interactive zero-knowledge proofs, which makes them a
natural building block for privacy-preserving applications such as group
signatures, anonymous and delegatable credentials, and privacy-preserving
payment systems. Many of these applications call for the signing process to
be distributed across several parties, so as to remove single points of
failure and limit the damage of key compromise, but efficient threshold
structure-preserving signatures had, until now, remained an elusive goal.

We gave two threshold structure-preserving signature schemes, built from the
structure-preserving signatures of Abe, Groth, Haralambiev, and Ohkubo and of
Kiltz, Pan, and Wee, together with generic MPC functionalities over both
field and group elements~\cite{C:AGHO11,C:KilPanWee15}. Both
constructions push the bulk of their computation into an offline
pre-processing phase, ahead of the message to be signed becoming known, so
that their online phase remains non-interactive and comparable in cost to
ordinary, non-threshold signing. We proved that both constructions achieve
adaptive TS-UF-1 security, inheriting the unforgeability of their underlying
structure-preserving signatures whenever the MPC functionalities they build
on are securely realised.

We evaluated both schemes using two realisations of these functionalities,
ATLAS in the honest-majority setting and MASCOT in the dishonest-majority
setting, and compared them against existing threshold structure-preserving
signatures. Our scheme built from the signature of Abe, Groth, Haralambiev, and Ohkubo
achieves optimally small threshold signatures of three group elements,
matching the size of its non-threshold counterpart, while our scheme built
from that of Kiltz, Pan, and Wee relies only on the standard SXDH
assumption~\cite{C:AGHO11,C:KilPanWee15}. In both cases, the offline phase
scales with the number of parties, while the online phase remains
constant-time and non-interactive, making the schemes practical for
high-throughput applications such as payments and transparency logs that
experience bursty traffic.

Beyond the constructions themselves, we gave the first multiparty secret
sharing of group elements, building on prior work extending secret sharing to
group elements, and implemented previously proposed threshold
structure-preserving signature schemes to allow for a concrete efficiency
comparison~\cite{NDSS:GPVRNC23}. Extending our approach to other structure-preserving
signatures, and to applications such as transparency logs and permissioned
ledgers discussed in Section~\ref{sec:eval}, would be valuable directions for
future work.
